% Fichier complété par tools/complete_missing_corrections.py.
% Source : TEC/TEC-05/19_Graham/19_Graham.tex
% Statut : rédigé (5 question(s)).
\begin{corrigebox}[Corrigé rédigé]
\CorrigeQuestion{1}
Le graphe comprend le bâti 4, l'arbre d'entrée 1, la roue mobile 2 et
                la sortie 3. Les contacts 2/4 et 2/3 sont des roulements sans glissement;
                2 est guidée radialement par rapport au porte-galet.
\CorrigeQuestion{2}
Au contact I, $V(I,2/4)=0$ donne
                $\omega(d/2-\lambda)=\omega_1 d/2$ (signes selon convention), donc
                $\omega=\omega_1 d/(d-2\lambda)$.
\CorrigeQuestion{3}
Au contact J, l'égalité des vitesses tangentielles donne
                $\omega(d_2/2+\lambda)=\omega_3d_3/2$ avec le signe d'engrènement.
\CorrigeQuestion{4}
En éliminant $\omega$, on obtient un rapport rationnel en $\lambda$ :
                $\omega_3/\omega_1=\pm[d(d_2+2\lambda)]/[d_3(d-2\lambda)]$.
\CorrigeQuestion{5}
On évalue cette expression pour $\lambda\in[12,23]$ mm et on trace
                une courbe monotone; les asymptotes éventuelles sont hors domaine si
                $d-2\lambda>0$.
\end{corrigebox}
